\paragraph{Proof strategy}
Theorem~\ref{thm:critical-limit}(i) gives \(\Gamma\)-convergence of the transport energy and convergence of the paired risk along bounded-energy plans; exact marginal recovery in part~(ii) makes the limiting kernel values asymptotically attainable in \(\Pi\). Part~(iii) gives uniform convergence of the optimized profiles on compact budget intervals, which together with the strict crossing in Theorem~\ref{thm:strict-crossing} locates the reversal budget \(s^*\).

\emph{Step 1: Localization and the energy lower bound.} Because the two ES thresholds are optimized independently, comparison of the objectives requires a common threshold window for stresses with bounded rescaled energy. In Appendix~\hyperref[proof:threshold-localization]{A.2}, Lemma~\ref{b-l2-all-optimal-thresholds-under-a-budget-of-order-varepsilon2} places every minimizing threshold in an \(O_M(\varepsilon)\) window. On bounded windows the lifted plans are locally compact; outside them, marginal conservation requires compensation for the projected local flux. The compensation bound above then gives the energy lower bound for arbitrary admissible plans. Choosing optimal plans yields the corresponding Wasserstein bound for their target laws.

\emph{Step 2: Convergence of the paired risk.} Bounded first moments make each hinge increment integrable, but need not give uniform integrability along convergent plans. The two objectives share the same large-displacement contribution, so we subtract \(\mathfrak q_0(0)\) from both. The relative increments are bounded on the common threshold window. Outside a normal window enlarged to account for the bounded reallocation \(H\), a nonzero relative increment requires a sufficiently large displacement, so the first moment controls the far-field contributions. The baseline integral is finite but need not converge. Since it translates both threshold objectives by the same amount independently of the threshold, it cancels exactly from their difference of minima.

\emph{Step 3: Exact marginal recovery.} Let \(\mathcal K_{\rm fin}^c\) denote probability kernels with compactly supported continuous weights on finitely many fixed displacement vectors and the remaining mass at zero, as in \eqref{eq:legacy-B.26}. For a fixed kernel in this class, local translations yield a preliminary law \(Q_\varepsilon\) with uniform marginal CDF expansion \(F_{i,\varepsilon}-F_i=-\varepsilon^2j_i+o(\varepsilon^2)\), where \(F_i\) is the reference CDF. The quantile maps \(F_i^{-1}\circ F_{i,\varepsilon}\) restore all three marginals at cost \(\varepsilon^2\sum_i\gamma_i\|j_i\|_1+o(\varepsilon^2)\), with \(O_\alpha(\varepsilon^2)\) displacements.

To see why this correction does not alter the leading paired risk, compare its displacement scale with the \(O_\alpha(\varepsilon)\) threshold band. Outside a sufficiently wide band, the paired hinge has gradient \(-\varepsilon h\) above both thresholds and zero below both, so multiplying by the restoration displacement gives an \(O_\alpha(\varepsilon^3)\) contribution. Inside the band the gradient is bounded, but the \(Q_\varepsilon\)-mass is only \(O_\alpha(\varepsilon)\), giving the same order after integration. The restoration error is therefore \(O_\alpha(\varepsilon^3)\), uniformly over all minimizing threshold pairs.

Approximation by \(\mathcal K_{\rm fin}^c\) in cost, paired risk, and the vague plan topology (Appendix A.4), followed by diagonal recovery, gives the general \(o(\varepsilon^2)\) error. Compactness, contraction toward the identity, and the common budget-Lipschitz bound then give locally uniform convergence of the optimized profiles.
